\documentclass[10pt,a4paper]{article}
\usepackage{fontspec}
\usepackage{polyglossia}
\setdefaultlanguage{french}
\setmainfont{Carlito}
\usepackage[top=1.2cm,bottom=1.2cm,left=1.4cm,right=1.4cm]{geometry}
\usepackage{microtype}
\usepackage{enumitem}
\usepackage{titlesec}
\usepackage{xcolor}
\usepackage{tabularx}
\usepackage[hidelinks]{hyperref}
\usepackage{graphicx}
\usepackage{tikz}
\usepackage{needspace}

% Polices d'icônes (Font Awesome, fichiers .ttf placés à côté de ce .tex)
\newfontfamily\FAsolid{fa-solid-900.ttf}[Path=./]
\newfontfamily\FAbrands{fa-brands-400.ttf}[Path=./]

\definecolor{accent}{HTML}{1E77B8}
\definecolor{muted}{HTML}{555555}
\definecolor{todo}{HTML}{C00000}

% ---------- PORTFOLIO (champ à renseigner) ----------
% 1) Remplace l'adresse ci-dessous (sans https://)  2) Passe \portfolioreadyfalse en \portfolioreadytrue
%    Tant que c'est « false », l'adresse s'affiche en rouge pour ne pas l'oublier.
\newif\ifportfolioready
\portfolioreadyfalse
\newcommand{\portfoliourl}{portfolio.example.com}

\pagestyle{empty}
\setlength{\parindent}{0pt}
\setlength{\parskip}{0pt}
\renewcommand{\arraystretch}{1.0}
\setlist[itemize]{leftmargin=1.1em,itemsep=1pt,topsep=1pt,parsep=0pt,label=\textcolor{accent}{\textbullet}}

\titleformat{\section}{\large\bfseries\color{accent}}{}{0pt}{}[\vspace{-5pt}\textcolor{accent}{\rule{\linewidth}{0.5pt}}]
\titlespacing*{\section}{0pt}{11pt}{5pt}

% Marqueur « à compléter » (rouge) : à supprimer avant envoi
\newcommand{\todo}[1]{\textcolor{todo}{[#1]}}

% \cvitem{date}{titre / lieu}{contenu}
% Ligne de contact : \contact{police}{code}{texte}
\newcommand{\contact}[3]{{\color{accent}\makebox[1.3em][c]{#1\symbol{#2}}}\,#3\par\vspace{3pt}}

% ---------- PHOTO ----------
% Dépose simplement ta photo dans le dossier du projet, sous l'un de ces noms :
%   img1.png, photo.jpg, photo.jpeg ou photo.png
% Elle est cadrée automatiquement en cercle (portrait, paysage ou carré, peu importe).
% Sans photo, un cercle gris « Photo » s'affiche.
\def\photofile{}
\IfFileExists{img1.png}{\def\photofile{img1.png}}{%
 \IfFileExists{photo.jpg}{\def\photofile{photo.jpg}}{%
 \IfFileExists{photo.jpeg}{\def\photofile{photo.jpeg}}{%
 \IfFileExists{photo.png}{\def\photofile{photo.png}}{}}}}

\newsavebox{\photobox}
\newlength{\photoshift}
\newcommand{\profilephoto}{%
  \ifx\photofile\empty
    \begin{tikzpicture}
      \fill[black!10] (0,0) circle (1.3cm);
      \node[muted] at (0,0) {\small Photo};
      \draw[accent,line width=0.8pt] (0,0) circle (1.3cm);
    \end{tikzpicture}%
  \else
    % on remplit le cercle : mise à l'échelle par la hauteur, ou par la largeur si l'image est en portrait
    \sbox{\photobox}{\includegraphics[height=2.6cm]{\photofile}}%
    \setlength{\photoshift}{0pt}%
    \ifdim\wd\photobox<2.6cm
      \sbox{\photobox}{\includegraphics[width=2.6cm]{\photofile}}%
      % portrait : on remonte un peu le cadrage vers le haut (là où se trouve en général le visage)
      \setlength{\photoshift}{\dimexpr(\ht\photobox-2.6cm)*3/10\relax}%
    \fi
    \begin{tikzpicture}
      \begin{scope}
        \clip (0,0) circle (1.3cm);
        \node[inner sep=0pt,yshift=-\photoshift] at (0,0) {\usebox{\photobox}};
      \end{scope}
      \draw[accent,line width=0.8pt] (0,0) circle (1.3cm);
    \end{tikzpicture}%
  \fi}

\newlength{\dcol}\setlength{\dcol}{2.2cm}
\newcommand{\cvitem}[3]{%
  \par\noindent\hangindent=\dcol\hangafter=1\raggedright
  \makebox[\dcol][l]{\color{muted}\small #1}#2\par
  \def\cvtmp{#3}\ifx\cvtmp\empty\else
    \noindent\hspace*{\dcol}\begin{minipage}[t]{\dimexpr\linewidth-\dcol\relax}\raggedright #3\end{minipage}\par
  \fi
  \vspace{6pt}}

% Icône ordinateur : lien vers le dépôt GitHub d'un projet
\newcommand{\repolink}[1]{\href{#1}{\color{accent}{\FAsolid\symbol{"F5FC}}}}
% Icône document : lien vers le rapport (Drive)
\newcommand{\reportlink}[1]{\href{#1}{\color{accent}{\FAsolid\symbol{"F15C}}}}
% Légende alignée à droite sur la ligne d'un titre de rubrique : \seclegend{code icône}{texte}
\newcommand{\seclegend}[2]{\hfill{\normalfont\small\color{muted}{\color{accent}\FAsolid\symbol{#1}}\,#2}}
\newlength{\scol}\setlength{\scol}{3.3cm}
\newcommand{\skill}[2]{\par\noindent\hangindent=\scol\hangafter=1\raggedright
  \makebox[\scol][l]{\bfseries\color{accent}#1}#2\par\vspace{2pt}}

\begin{document}
\raggedright

% ---------- EN-TÊTE ----------
\noindent
\begin{minipage}[c]{2.7cm}\profilephoto\end{minipage}\hspace{4mm}%
\begin{minipage}[c]{\dimexpr\linewidth-2.7cm-4.8cm-8mm\relax}\raggedright
  {\fontsize{16.5}{20}\selectfont\bfseries\color{accent}\mbox{Kayossi Olouwa Charmant BALOGOUN}}\\[6pt]
  {\fontsize{13}{16}\selectfont\bfseries Modélisation et Simulation Numérique}\\[2pt]
  {\color{muted} Élève ingénieur en Génie Mathématique et Modélisation}
\end{minipage}\hfill
\begin{minipage}[c]{4.8cm}\raggedright\small
  \contact{\FAsolid}{"F0AC}{\href{https://charmantkoc.github.io/}{https://charmantkoc.github.io/}}
  \contact{\FAsolid}{"F0E0}{\href{mailto:charmantbalogoun@gmail.com}{charmantbalogoun@gmail.com}}
  \contact{\FAbrands}{"F09B}{\href{https://github.com/CharmantKOC}{CharmantKOC}}
  \contact{\FAbrands}{"F08C}{\href{https://www.linkedin.com/in/charmantblg}{charmantblg}}
\end{minipage}

\vspace{2pt}

% ---------- PROFIL ----------
\section*{Profil}
Élève ingénieur en fin de cycle, spécialisé en modélisation et simulation numérique, avec un socle solide en
analyse numérique et en équations différentielles. Mes intérêts portent sur la modélisation des systèmes dynamiques,
la simulation de phénomènes complexes et ses interactions avec l'intelligence artificielle (prédiction, optimisation,
aide à la décision). Je recherche un stage de fin d'études ou une poursuite en master ou en doctorat,
dans un environnement de recherche et d'innovation en modélisation, contrôle et calcul scientifique, sur des problématiques réelles.

% ---------- FORMATION ----------
\section*{Formation}
\cvitem{2024-Présent}
 {\textbf{Cycle ingénieur}, Modélisation et Simulation Numérique, École Nationale Supérieure de Génie Mathématique et Modélisation (ENSGMM), UNSTIM, Abomey, Bénin}
 {\small Cours majeurs : modélisation mathématique, analyse et simulation numérique, optimisation, calcul scientifique, simulation stochastique,
  modélisation statistique, génie logiciel.}
\cvitem{2022-2024}
 {\textbf{Classes préparatoires}, Sciences de l'ingénieur, Institut National Supérieur des classes Préparatoires aux Études d'Ingénieur (INSPEI), UNSTIM, Abomey, Bénin \hfill \textit{Mention Très Bien}}
 {\small Cours majeurs : analyse, algèbre, analyse complexe, analyse numérique, graphes et optimisation, recherche opérationnelle,
  modélisation des phénomènes aléatoires, algorithmique, C, MATLAB.}
\cvitem{2022}
 {\textbf{Baccalauréat série C}, Sciences techniques, Lycée Houffon d'Abomey, Bénin \hfill \textit{Mention Très Bien}}{}

% ---------- ÉCOLES ----------
\section*{Écoles et formations}
\cvitem{Oct. 2026}
 {\textbf{École de recherche CIMPA « Mathématiques appliquées et Intelligence artificielle »}, IMSP, Dangbo, Bénin}
 {\small Cours : apprentissage automatique et théorie du contrôle (D.\ Prandi, CNRS) ; mathématiques numériques et apprentissage
  des EDP aux signaux cérébraux (P.\ Minakowski) ; optimisation sur les graphes ; calcul parallèle et HPC.}
\cvitem{2024-2025}
 {\textbf{Certificat en Intelligence artificielle}, programme FORCE-N, Université Numérique Cheikh Hamidou Kane - Mastercard Foundation}{}
\cvitem{Août 2022}
 {\textbf{École d'été sur l'Intelligence artificielle (EEIA 2022)}, Fondation Vallet - Bénin Excellence - PNUD, Godomey}{}

% ---------- PROJETS ----------
\section*{Projets\seclegend{"F5FC}{Dépôt GitHub}}
\cvitem{2026}
 {\textbf{Modélisation de la décomposition spinodale par l'équation de Cahn-Hilliard} \hfill \repolink{https://github.com/CharmantKOC/Modelisation-et-simulation-numerique-de-la-decomposition-spinodale-par-equation-de-Cahn-Hilliard}}
 {\small Modélisation et simulation numérique d'une EDP non linéaire d'ordre quatre, avec discrétisation spectrale par transformée en cosinus discrète (Python).}
\cvitem{2026}
 {\textbf{Application des Physics Informed Neural Networks} \hfill \repolink{https://github.com/CharmantKOC/Physics-Informed-Neural-Networks-PINNs-resolution-d-equations-differentielles}}
 {\small Résolution de trois équations différentielles par réseaux de neurones: projectile, oscillateur amorti et modèle SIR (Python).}
 \cvitem{2025}
 {\textbf{Motifs de Turing : modèle de réaction-diffusion de Gray-Scott} \hfill \repolink{https://drive.google.com/file/d/1zA777E19PlPkxRNqL2tHVwHHl5-p4-MG/view?usp=sharing}}
 {\small Simulation par différences finies du système de Gray-Scott et analyse numérique de la formation de motifs (Python).}
\cvitem{2025}
 {\textbf{Reconnaissance vocale de la langue Fon par chaînes de Markov}, ENSGMM \hfill \repolink{https://github.com/CharmantKOC/Markov_Fon}}
 {
   \small Modélisation par chaînes de Markov, visualisation des modèles (NetworkX) et interface Streamlit (Python, Matplotlib).
  }
\cvitem{2024}
 {\textbf{Analyse de commentaires clients}, ENSGMM}
 {\small Traitement de commentaires et tableau de bord interactif (Python, Pandas, WordCloud, Streamlit).}

% ---------- EXPÉRIENCE ----------
\section*{Expérience professionnelle\seclegend{"F15C}{Rapport}}
\cvitem{Juil.-sept. 2026}
 {\textbf{Stage académique}, Systèmes de données et cartographie numérique, Cabinet S Lab, Abomey-Calavi \hfill \reportlink{https://drive.google.com/file/d/1_ELQTKoKy-eciZEjD1EQv7nKonXfXU5Y/view?usp=sharing}}
 {\small Systèmes d'information géographique (SIG), traitement de données spatiales, cartographie par drone et développement d'une plateforme numérique métier (Python, Django, Leaflet.js).}
\cvitem{Août 2025}
 {\textbf{Stage académique}, Statistiques et analyse de données, Direction Départementale de la Santé du Zou, Abomey \hfill \reportlink{https://drive.google.com/file/d/1j5Kd7Ho0rU9hWMbA87faC99PteDjQBc5/view?usp=sharing}}
 {
   \small Analyse et modélisation des tendances épidémiologiques du paludisme dans le département du Zou ; tableau de bord interactif pour le suivi et l'aide à la décision (Python, Matplotlib, Streamlit).
  }

% ---------- CERTIFICATIONS ----------
\section*{Certifications et conférences}
\cvitem{Juin 2026}{\small \textbf{IOCMA 2026}, 2\textsuperscript{nd} International Online Conference on Mathematics and Applications, MDPI}{}
\cvitem{Sept. 2026}{\small \textbf{ML Summer School 2026}, Cohere Labs Open Science Community}{}
\cvitem{Déc. 2025}{\small \textbf{Modélisation par équations structurelles sous R} (santé, environnement, agronomie), LABEF, Université d'Abomey-Calavi}{}
\cvitem{Juin 2024}{\small \textbf{Mise en œuvre des réseaux de neurones génératifs profonds} (usurpation d'identité numérique), ANSALB / CITA, Cotonou}{}
%\cvitem{}{\small Certificats Kaggle}{}

% ---------- DISTINCTIONS ----------
\Needspace{12\baselineskip}
\section*{Distinctions : hackathons et concours}
\cvitem{2026}
 {\textbf{1\textsuperscript{er} prix}, Hackathon Confédération des Institutions Financières
d'Afrique de l'Ouest, DigiCoop-WA+, équipe Vigie-Coop}
 {\small  Conception d'un système de filtrage des clients et détection d'opérations frauduleuses dans les coopératives financières.}
\cvitem{2025}
 {\textbf{3\textsuperscript{e} prix national}, NASA International Space Apps Challenge, équipe CMAK}
 {\small Prévision de la qualité de l'air à partir de données d'observation de la Terre (machine learning, API FastAPI).}
\cvitem{2025}
 {\textbf{3\textsuperscript{e} prix}, FRIARE Hackathon « Smart Cities », équipe CCP}
 {\small Détection automatique de plaques d'immatriculation via caméras routières, vérification en base de données et alertes de non-conformité (Python, OpenCV, YOLOv8, Pandas).}
\cvitem{2024}
 {\textbf{Finaliste}, Kyfax AI Contest, ENSGMM, équipe JoLyCH}
 {\small Agent IA jouant au jeu traditionnel Sholo Guti (16 soldats), avec optimisation des décisions (Python).}
\cvitem{2022}
 {\textbf{3\textsuperscript{e} prix}, MIFY Artificial Intelligence Contest, équipe Yesunana}
 {\small Agent IA jouant au Morabaraba.}
%\cvitem{2025}
 %{\textbf{Participations}}
 %{\small Coupe d'Afrique des Nations en Science des Données (DataTour 2025), phase nationale (équipe ChAy).}

% ---------- OUTILS ----------
\section*{Outils et technologies}
\skill{Programmation}{C, MATLAB, Python, Fortran, R}
\skill{Données et IA}{NumPy, SciPy, Pandas, Matplotlib}
\skill{Développement}{FastAPI, Streamlit}
\skill{Environnement}{Linux, Git, LaTeX}
%\skill{CAO}{AutoCAD, TopSolid}

% ---------- LANGUES ----------
\section*{Langues et centres d'intérêt}
\skill{Langues}{Français : C2 \,·\, Anglais : B1}
\skill{Centres d'intérêt}{Basketball, piano}

\end{document}